\subsection{Eight build algorithms}
\label{subsec:eight-build-algorithms}

The seventh entry of the Dirac--Igusa pro-object is the explicit
construction calculus: eight finite-stage algorithms which, run in
sequence, take the chosen Bridgeland heart on
$D^b(\operatorname{Coh}(K3\times E))$ and the imported BKM target
$\mathfrak g_{\Delta_5}$ as inputs and produce the pro-object
$\mathfrak D_X^{\mathrm{DI}}$.  Each algorithm has a typed input,
a typed output, and a finite obstruction tuple; the pro-object is the
inverse limit when every obstruction tuple vanishes compatibly with
HN transitions.  The full pseudocode and per-algorithm
input/output/obstruction tables are
Appendix~\ref{app:eight-build-algorithms}.

The discipline is finite-first.  Every operation acts on
finite-type stacks and finite-dimensional Borel--Moore chains at
fixed HN height; the inverse limit is taken only after every finite
stage is supplied with compatible transition maps.  The eight
algorithms are mutually dependent only in the order listed.  For
$R\in\mathcal R_{\mathrm{HN}}$, the algorithms assume:
$[\textup{K3$\times$E-fin}]$ (finite-type semistable substacks
at bounded HN height),
$[\textup{K3$\times$E-atlas}]$ (finite d-critical/cosection
Hall atlas with Joyce--Upmeier orientation lines), and
$[\textup{ML}]$ (Mittag--Leffler descent).

\subsubsection*{Algorithm A1: charge window}

\textit{Input.} Stability sector $(\sigma,S)$, finite charge window
$\Gamma_R^{\mathrm{test}}\subset\Gamma_X^{\mathrm{form}}$.

\textit{Output.} Finite saturated downward window
$F_{\sigma,S}(R)\subset\Gamma_X^{\mathrm{HN}}(R)$, the dyonic
formal Mukai/effective-charge support, the Gram map
$\Pi_X:\Gamma_X^{\mathrm{form}}\to\Gamma_{\mathrm{gram}}$, the cocycle
$B$, and the normal-ordered extension
$\widehat\Gamma_R=\Gamma_X^{\mathrm{form}}\oplus_B\Gamma_{\mathrm{gram}}$.

\textit{Procedure.}
\begin{enumerate}
\item Form the Liu charge $L_F(n)=Z_S(Rp_{S,*}(F\otimes p_E^*H^n))$
\cite{LiuProductStability} for $\sigma_S$ a rational K3 stability
condition and a degree-one line bundle $H$ on $E$.  Restrict to
charges with bounded $h_S$-height $\le R$; the resulting set is
$F_{\sigma,S}(R)$.
\item Compute the Gram map and cocycle $B(c,c')$ on
$F_{\sigma,S}(R)\times F_{\sigma,S}(R)$ via the Mukai dictionary of
\S\ref{subsec:d6-d2-d0-dictionary-appendixC}; check the symmetric
bilinear $2$-cocycle condition.
\item Form $\widehat\Gamma_R$ as the central extension by
$B$-cocycle.
\end{enumerate}

\textit{Finite obstructions.}
$[\textup{K3$\times$E-fin}]$ on the Liu charge bound; $B$
cocycle property.

\subsubsection*{Algorithm A2: finite Hall stage and atlas}

\textit{Input.} Output of A1, plus the cosection-twisted reduced
d-critical heart of \S\ref{subsec:cosection-twisted-orientation}.

\textit{Output.} The finite Hall stage
$\mathcal C_{\sigma,S,\le R}$, finite-type derived Artin enhancements
$\mathfrak M_{\widehat c,R},
\mathfrak E_{\widehat c,\widehat c',R},
\mathfrak F^{(2)}_{\widehat c,\widehat c',\widehat c'',R}$ of object,
extension, and two-step flag stacks; reduced cosections; reduced
vanishing cycles $\Phi^{\mathrm{red}}$; orientation lines
$o^{\mathrm{red}}$.

\textit{Procedure.}
\begin{enumerate}
\item Restrict the Bridgeland heart to charges in $F_{\sigma,S}(R)$;
form the finite-type derived enhancements with PTVV $(-1)$-shifted
symplectic structures \cite{PTVV2013} on the truncations; descend to
Joyce d-critical structures \cite{BBDJS2015}.
\item Compute the K3 semiregularity cosection
$\operatorname{CS}_{\mathcal C}$ on each retained charge stratum;
verify surjectivity on reduced strata and filtered Hall additivity
on extension stacks.
\item Build Joyce--Upmeier orientation lines \cite{JoyceUpmeier} with
chosen square roots, multiplicative under reduced extension
correspondences.
\end{enumerate}

\textit{Finite obstructions.}
$[\textup{K3$\times$E-atlas}]$; finite residual inertia after
scalar, $E$-translation, and wrapped rigidifications.

\subsubsection*{Algorithm A3: hybrid wrapped factorization datum}

\textit{Input.} Output of A2.

\textit{Output.} The hybrid factorization carrier
$\mathcal F_{X,\sigma,S,\le R}^{\mathrm{hyb}}$ of
\S\ref{subsec:hybrid-ran-carrier}: rigidified wrapped prequotients
$\mathcal M_{\eta,R}^{\mathrm{wr,rig}}$ with anchors
$\lambda_{\eta,R}$; closed-configuration local sectors
$\mathcal F_R^{\mathrm{loc}}$; wrapped sectors
$\mathcal G_R^{\mathrm{wr}}$; the eight-word two-step binary flag
atlas; the higher coloured residual
$\mathfrak o^{\mathrm{col}}_R$; quotient functor $Q_{E,R}$ and
comparison isomorphisms $\theta^Q_{\mu,R}$; protected integration
$I_R^{\mathrm{prot}}$.

\textit{Procedure.}
\begin{enumerate}
\item Split $\Gamma_R$ into $\Gamma_R^{\mathrm{loc}}$
($b_R^{\mathrm{geom}}=0$) and $\Gamma_R^{\mathrm{wr}}$
($b_R^{\mathrm{geom}}>0$).
\item For $\alpha\in\Gamma_R^{\mathrm{loc}}$: form the
closed-configuration incidence stacks
$\mathfrak M_{\alpha,R,I}^{\mathrm{loc,cl}}$ and the configuration
maps $\pi_{\alpha,R,I}:\mathfrak M_{\alpha,R,I}^{\mathrm{loc,cl}}
\to E^I$, with vanishing-cycle/orientation pushforwards
$\mathscr H_{\alpha,R,I}^{\mathrm{loc}}$.  Define
$\mathcal F_{\alpha,R}^{\mathrm{loc}}(U)$ by symmetric-group descent
on $E^I$.
\item For $\eta\in\Gamma_R^{\mathrm{wr}}$: form the rigidified
prequotient $\mathcal M_{\eta,R}^{\mathrm{wr,rig}}\to
\mathcal M_{\eta,R}^{\mathrm{wr}}$; choose anchor
$\lambda_{\eta,R}$ with translation weight one; verify the typed
finite anchor residual $\mathfrak o^\lambda_R=0$ on each component.
\item Build the local-local, ordered local-wrapped, and
wrapped-wrapped extension correspondences with closed-support
finite-type witnesses, exceptional-pushforward admissibility, and
Thom--Sebastiani identifications.
\item Build the eight-word two-step binary flag stacks
$\mathfrak{Flag}^w_{\xi_1,\xi_2,\xi_3,R}$ for
$w\in\{\mathrm L,\mathrm W\}^3$ and verify
$\mathfrak o^{\mathrm{assoc}}_{w,R}=0$ on the four-input pentagon.
\item Build the quotient pseudofunctor
$\mathsf{Corr}^{E,\mathrm{hyb}}_R\to
\mathsf{Corr}^{\mathrm{red},\mathrm{hyb}}_R$ and the comparison
isomorphisms $\theta^Q_{\mu,R}$; verify the four quotient residuals.
\item Build $I_R^{\mathrm{prot}}$ on the reduced functor; verify the
five integration defects vanish and the Gram-trace exponent law
$I_R^{\mathrm{prot}}(Q_{E,R}x)\in s^{m_R(\gamma)}\C[q^{\pm1},r^{\pm1}]$
with $m_R(\gamma)=\operatorname{pr}_3\overline\Pi_X(\gamma)$.
\end{enumerate}

\textit{Finite obstructions.}  $\mathfrak O_{\mathrm{hyb},R}$
(\S\ref{subsec:hybrid-ran-carrier}).

\subsubsection*{Algorithm A4: orientation descent and Weyl transport}

\textit{Input.} Output of A3.

\textit{Output.} The (O1) and (O1)$^+$ data: null-trivializations of
the reduced and equivariant orientation gerbes,
$E[N]$-linearizations with vanishing degree-one characters, and
Weyl-equivariant lifts $\tau_i$ ($i=1,2,3$) compatible with the
quotient-orientation package $\mathfrak q_{\mathcal C}$.

\textit{Procedure.}
\begin{enumerate}
\item For each charge stratum $\mathcal C$: compute the reduced
gerbe $\alpha^{\mathrm{red}}_{\mathcal C}$, free $E$-class
$\alpha^{E,\mathrm{free}}_{\mathcal C}$, and equivariant
finite-stabilizer classes $\widetilde\beta^{H}_{\mathcal C,S}$ for
the actual stabilizers $H\subset E_{\mathrm{tors}}$ on
finite-stabilizer strata.
\item Verify the Klein-four condition at $H=E[2]$:
$(b_{20},b_{11},b_{02})=(r_1,r_1+r_2+r_3,r_2)$ with
$r_1=r_2=r_3=0$; for full-level even $N\ge 4$, verify the
two-primary residual on $E[N]_{(2)}$ including the mixed coefficient
$A_{12}^{(N)}$.
\item Choose zero $H^1(BH;\mathbb F_2)$-linearization characters
$\lambda^H_{\mathcal C,S}=0$.
\item For each simple type-II reflection $s_{\delta_i}$, supply the
lift $\tau_i:o_{\mathcal C}\to s_{\delta_i}^* o_{\mathcal C}$ on the
finite partial action groupoid $\mathcal G_R$; compute the
$2$-cocycle $c_{o,R}\in Z^2(\mathcal G_R;\underline{\mathbb F}_2)$
and verify $[c_{o,R}]=0$ with chosen $1$-cochain.
\item Verify $\tau_i^2=1$ and the Coxeter relations on the
$W^{(2)}(\Lambda_{II}^{2,1})$-orbit; verify
$\omega_{i,\mathcal C}=0$ for each torsor defect of the
quotient-orientation transport.
\end{enumerate}

\textit{Finite obstructions.}  $\mathfrak O_{\mathrm{or},R}$
(\S\ref{subsec:cosection-twisted-orientation}); the (O2) wall-atlas
is built next.

\subsubsection*{Algorithm A5: Hall product, coproduct, and primitive
projection}

\textit{Input.} Outputs of A2--A4.

\textit{Output.} The compact Hall product $m_R^{\mathrm{red}}$,
collision coproduct $\Delta_R^{\mathrm{ch}}$, primitive projection
$P_R^{\Prim}$, and Hopf pairing $\langle\,,\,\rangle_R$ on the
cosection-reduced d-critical heart, with normal-ordered Gram descent
$\overline\Pi_{X,*}^\Theta$.

\textit{Procedure.}
\begin{enumerate}
\item Form $m_R^{\mathrm{red}}=q_!\circ\operatorname{TS}^{\mathrm{red}}_R
\circ p^*$ on each binary correspondence (LL, LW, WL, WW); verify
two-step pentagon coherence on the eight-word atlas.
\item Form $\Delta_R^{\mathrm{ch}}$ by pull-push on the $2$-step flag
stack; verify coassociativity and counitality after $Q_{E,R}$.
\item Form $\langle\,,\,\rangle_R$ on the primitive layer; verify
Hopf adjointness with the coproduct on the radical quotient.
\item Compute the Frobenius defect $\mathfrak o_{F,R}$ via the
CY-3 Serre-functor identity \cite{KSCoHA}; verify
$\mathfrak o_{F,R}=0$.
\item Compute the radical $\operatorname{Rad}_R^\Pi$; verify
Hopf-coideal and Lie-ideal stability.
\item Form the chain-level normal-ordered Gram descent
$\overline\Pi_{X,*}^\Theta$ via the negative-cyclic
Loday--Quillen--Tsygan route or the PTVV--Joyce orientation route
\cite{LodayQuillen1984,Tsygan1983,LodayCyclicHomology,PTVV2013};
verify the triple homotopy compatibility
$d_{\mathrm{Hoch}}\Theta_\Pi=B_{\mathrm{ch}}$ and the
$\widehat\Gamma_R$-grading absorption.
\end{enumerate}

\textit{Finite obstructions.}  $\mathfrak O_{\mathrm{Hall},R}$,
including $\mathfrak o_{F,R}$, the Hopf-coideal residual, and the
$\Theta$-pentagon residual.

\subsubsection*{Algorithm A6: Dirac block, Pfaffian line, type-II
wall atlas}

\textit{Input.} Outputs of A4 and A5.

\textit{Output.} The first-order Dirac block decomposition
$\mathfrak D_{X,R}^{D_0}=\bigoplus_\gamma J_\gamma$, the Pfaffian
line $\mathscr L_{\mathrm{Pf},R}$ with squaring isomorphism, the
Pfaffian-to-automorphic comparison $\iota_{\mathrm{aut},R}$, the
normalized section $\operatorname{pf}_{X,R}$, and the type-II wall
atlas (rank-one normal model on every retained component, boundary
stratum, and overlap).

\textit{Procedure.}
\begin{enumerate}
\item Form the Pfaffian line
$\mathscr L_{\mathrm{Pf},R}=\bigotimes_{\mathcal C}
\operatorname{Pf}(K^{\mathrm{red}}_{\mathcal C,R})$ on the
cosection-reduced quotient \cite{BBDJS2015}; verify the squaring
isomorphism.
\item Form $J_\gamma$ on
$P^{\Pi,+}_{R,\gamma}\oplus(P^{\Pi,+}_{R,\gamma})^\vee[1]$ for each
$\gamma\in\Gamma_R^{\Pi,+}\setminus\{0\}$; verify Pfaffian
contribution $(1-x^\gamma)^{\sdim P^{\Pi,+}_{R,\gamma}}$.
\item Verify the support, parity, and leading-coefficient residuals
$\mathfrak o^{\mathrm{supp}}_{R,\gamma}=
\mathfrak o^{\mathrm{par}}_{R,\gamma}=
\mathfrak o^{\mathrm{lead}}_R=0$ on the test window.
\item For each simple type-II wall $\Hpl_{\delta_i}$ ($i=1,2,3$):
build the rank-one Pfaffian crossing chart on every retained
component and boundary stratum; verify
$N_{\delta_i}^{\mathrm{Pf}}=1$ and the local sign
$\varepsilon^{\mathrm{loc}}_{\delta_i,R}=-1$; verify chart-overlap
agreement of the normal Pfaffian units, tangent/normal splittings,
quotient orientations, and protected integration maps; verify the
half-Hilbert orbit sign sum (size $4096=2^{12}$).
\item Supply the Pfaffian-to-automorphic comparison
$\iota_{\mathrm{aut},R}:\mathscr L_{\mathrm{Pf},R}\to\mathcal L^5
\otimes\nu_{\Delta_5}$ compatible with cusp trivialization and
squaring.
\end{enumerate}

\textit{Finite obstructions.}  (D0-Pf), (O2)/hybrid compatibility
residuals, and the Pfaffian-to-automorphic comparison residuals.

\subsubsection*{Algorithm A7: Koszul source coalgebra and cobar
comparison}

\textit{Input.} Outputs of A3 and A5.

\textit{Output.} The conilpotent chiral source coalgebra $C_{X,R}$,
finite filtration $F^{\mathrm{co}}_{R,\bullet}$, collision coproduct
$\Delta_R^{\mathrm{ch}}$, bar comparison $b_{X,R}$, and cobar
comparison $\Theta_{\mathrm{Kos},R}$.

\textit{Procedure.}
\begin{enumerate}
\item Choose an augmentation
$\varepsilon_{X,R}^{\mathrm{hyb}}:\mathcal F_{X,\sigma,S,\le R}^{\mathrm{hyb}}
\to k\mathbf 1$ and verify bounded length $N_R$.
\item Form the canonical source-bar coalgebra
$C^{\mathrm{bar}}_{X,R}=\bar B_E^{\mathrm{ch}}(\mathcal F_{X,\sigma,S,\le R}^{\mathrm{hyb}})$
\cite[\S3.7]{BD} with bar-length filtration and deconcatenation
collision coproduct.
\item Verify the eight-word coassociativity, bar-length conilpotence,
and transition compatibility; set $b_{X,R}=\mathrm{id}$.
\item Verify the source-internal counit
$\varepsilon^{\mathrm{src}}_{\mathrm{bar},R}:\Omega_E^{\mathrm{ch}}
C^{\mathrm{bar}}_{X,R}\to\mathcal F_{X,\sigma,S,\le R}^{\mathrm{hyb}}$
is a quasi-isomorphism in the chosen completed/coderived category
\cite[Theorem~6.2.2]{FrancisGaitsgoryCKD}.
\item Construct the source-to-target chiral algebra
quasi-isomorphism $\vartheta_R:
\mathcal F_{X,\sigma,S,\le R}^{\mathrm{hyb}}\xrightarrow{\simeq}
\mathsf{Den}_{\Delta_5,E,\le R}$ preserving differential, product,
unit, transition maps, and primitive restriction.
\item Set $\Theta_{\mathrm{Kos},R}=
\vartheta_R\circ\varepsilon^{\mathrm{src}}_{\mathrm{bar},R}$;
verify the Koszul Mittag--Leffler condition on the cones of $b$ and
$\Theta_{\mathrm{Kos}}$.
\end{enumerate}

\textit{Finite obstructions.}  $\mathfrak O_{\mathrm{Kos},R}$,
including $\mathfrak o^\Omega_R=
(\mathfrak o^{\varepsilon,\mathrm{src}}_R,
\mathfrak o^\vartheta_R)$.

\subsubsection*{Algorithm A8: primitive recognition (cofinal
finite-window)}

\textit{Input.} Outputs of A5--A7, plus the imported Borcherds--Kac
target datum $\mathscr B_{\Delta_5}$.

\textit{Output.} For every cofinal saturated finite window $W_\nu$,
the recognition data (R0)--(R8); inverse-limit isomorphism
$P_X^{\Pi,\mathrm{red}}\cong\mathfrak g_{\Delta_5}$ and chiral
identification $\Omega_E^{\mathrm{ch}}C_X\simeq\mathsf{Den}_{\Delta_5,E}$.

\textit{Procedure.}
\begin{enumerate}
\item For each $\nu$, identify the saturated downward window $W_\nu$
in $\mathcal R_+$.
\item Build the simple-root representatives $e_i^X,f_i^X,h_i^X$ in
$P^{\Pi,\mathrm{red}}_{X,\pm\alpha_i}$ for $\alpha_i\in W_\nu$; verify
parity assignments via the GN multiplicity fibres.
\item Verify the Borcherds--Kac relations (Cartan, Chevalley, real
Serre, orthogonality, super sign) on every relation whose source and
target degrees lie in $W_\nu\cup(-W_\nu)\cup\{0\}$.
\item Compute the full finite parity dimensions
$\dim P^{\Pi,\mathrm{red}}_{X,\beta,\overline p}$ ($p=0,1$) for
$\beta\in W_\nu$; verify two-integer equality with
$\rootmult_{\overline p}^{\mathrm{GN}}(\beta)$.
\item Compute the Hopf pairing matrices in parity blocks; verify
kernel agreement with GN/Kac, coideal stability, and transition
compatibility.
\item Verify the no-extra-relations equality
$\ker\pi_\nu=
(\overline{J_{\mathrm{BK}}+\operatorname{Rad}_{\mathrm{GN}}})_{\le W_\nu}$.
\item Verify generation by Cartan and simple representatives;
compute PBW associated gradeds and verify strict transition
preservation.
\item Verify exact triangular completion: $\lim^1=0$ on the
triangular finite-window systems.
\end{enumerate}

\textit{Finite obstructions.}  $\mathfrak O_{\mathrm{Rec},\nu}$
(\S\ref{subsec:primitive-recognition}); for the chiral
identification, the inverse-limit Koszul Mittag--Leffler condition
of \S\ref{subsec:koszul-source}.

\subsubsection*{Construction summary}

The eight-step calculus reads
\[
\begin{array}{rcl}
(\sigma,S,R) & \xrightarrow{\textrm{A1}} &
F_{\sigma,S}(R),\widehat\Gamma_R\\
& \xrightarrow{\textrm{A2}} &
\mathcal C_{\sigma,S,\le R},\Phi^{\mathrm{red}},o^{\mathrm{red}}\\
& \xrightarrow{\textrm{A3}} &
\mathcal F_{X,\sigma,S,\le R}^{\mathrm{hyb}}\\
& \xrightarrow{\textrm{A4}} &
\textup{(O1)},\textup{(O1)}^+\\
& \xrightarrow{\textrm{A5}} &
m_R^{\mathrm{red}},\Delta_R^{\mathrm{ch}},
\langle\,,\,\rangle_R,\overline\Pi_{X,*}^\Theta\\
& \xrightarrow{\textrm{A6}} &
\mathfrak D_{X,R}^{D_0},\mathscr L_{\mathrm{Pf},R},
\operatorname{pf}_{X,R},\textup{(O2)}\\
& \xrightarrow{\textrm{A7}} &
C_{X,R},b_{X,R},\Theta_{\mathrm{Kos},R}\\
& \xrightarrow{\textrm{A8}} &
\textup{(R0)\text{--}(R8)\ on\ }W_\nu
\end{array}
\]
The pro-object
\[
\mathfrak D_X^{\mathrm{DI}}
=\varprojlim_R\bigl(
\mathcal A_{X,R}^{E_3},
F_{X,R}^{\mathrm{hyb}},
\Gamma_{X,R},
\Pi_X,
\Phi_R,
o_R,
H_R,
P_R^\Pi,
C_{X,R},
\Theta_{\mathrm{Kos},R},
\mathscr L_{\mathrm{Pf},R},
\operatorname{pf}_{X,R},
\varepsilon_{o,R},
\mathrm{Rec}_R\bigr)
\]
is the inverse limit when every finite obstruction tuple vanishes
compatibly.

\subsubsection*{Theorem statements at finite stage}

\begin{theorem}[Finite-stage Dirac--Igusa pro-object;
\textit{conditional on the finite obstruction tuples}]
\label{thm:finite-stage-dirac-igusa-pro-object}
Suppose Algorithms A1--A8 are run at heights
$R\in\mathcal R_{\mathrm{HN}}$ with all finite obstruction tuples
vanishing and HN transitions strict.  Then the eight outputs assemble
into a finite-stage Dirac--Igusa datum
$\mathfrak D_{X,R}^{\mathrm{DI}}$ at every $R$, with strict
transition maps to $\mathfrak D_{X,R'}^{\mathrm{DI}}$ for $R'\le R$,
and the inverse limit $\mathfrak D_X^{\mathrm{DI}}=\varprojlim_R
\mathfrak D_{X,R}^{\mathrm{DI}}$ exists in the completed
chiral/Hall framework.
\end{theorem}

\begin{theorem}[Pfaffian--Dirac identification at finite stage;
\textit{conditional on (D0), (O1), (O1)$^+$, (O2) at $R$}]
\label{thm:pfaffian-dirac-finite-stage}
Suppose the finite stage outputs of A1--A8 vanish on every relevant
obstruction tuple at height $R$.  Then on the cusp expansion
\[
\operatorname{pf}_{X,R}|_{\mathfrak c_\infty}
=64q^{1/2}r^{1/2}s^{1/2}\prod_{\gamma\in\Gamma_R^{\Pi,+}\setminus\{0\}}
(1-x^\gamma)^{\sdim P^{\Pi,+}_{R,\gamma}},
\]
the Pfaffian section matches the Borcherds product on the test
window, with the orientation character
$\varepsilon_{o,R}|_{W^{(2),\le R}}=\nu_{\Delta_5}|_{W^{(2),\le R}}$
on the type-II reflection truncation.  The recognition map
$\iota_R^{\Prim}:H^0\Prim_{\mathrm{prot}}(\overline\Pi_{X,*}^\Theta
\mathcal F_{X,\sigma,S,\le R}^{\mathrm{hyb}})\to\mathfrak g_{\Delta_5,\le R}$
is an isomorphism on the saturated finite window.
\end{theorem}

\subsubsection*{Status under the cross-volume probes}

Probes A005, A011, A026, A037, A042 verify components of A2--A4 at
the test window; probes A074, A076, A082 verify components of A4 at
the (O1) descent; A047, A049, A089 verify components of A4--A6 at
(O1)$^+$; A047, A092, A100 verify components of A6 at (O2).  The
A8 cofinal recognition ledger is the recognition probes
A049--A089.  Discharge of A1--A6 at the relation-closed window
beyond $W_{\le 3}$ is the open frontier of the cross-volume Master
Critique.

Algorithms A1--A8 run cleanly only when four obstruction classes
vanish.  Entry $(\mathrm L 8)$ names these obstructions $(D0)$,
$(O1)$, $(O1)^+$, $(O2)$, locates each in the construction, and
tracks the discharges through
Appendix~\ref{app:obstruction-discharges} in
\S\ref{subsec:four-named-obstructions-ch5}.
